\documentclass{article}

\usepackage[urw-garamond]{mathdesign}

\usepackage{geometry}
\usepackage{verbatim}
\usepackage{cite}
\usepackage{listings}

\lstdefinelanguage{siexml}
  {morekeywords={xml,sie,version,xmlns,decoder,loop,var,bits,type,endian,
                 value,group,tag,id,format,read,sample,octets,encoding,
                 ch,test,index,dim,set,data,name},
   sensitive=false,
   morecomment=[s]{<!--}{-->},
   morestring=[b]",
   showstringspaces=false,
   columns=fullflexible
  }

\lstdefinelanguage{rng}{
  morekeywords={attribute,div,element,
    empty,external,grammar,include,inherit,list,
    mixed,notAllowed,parent,string,
    text,token},
  sensitive=false,
  morecomment=[l]{\#},
  morestring=[b]",
  columns=fullflexible
}

\lstnewenvironment{grammar}
  {\vbox\bgroup\lstset{language=rng,frame=simple}}
  {\egroup}

\lstnewenvironment{siexmlexample}
  {\vbox\bgroup\lstset{language=siexml,frame=simple}}
  {\egroup}

\newcommand{\relaxng}[0]{\hbox{RELAX~NG}}
\newcommand{\xml}[1]{\texttt{#1}}
\newcommand{\var}[1]{\texttt{#1}}
\newcommand{\lit}[1]{\texttt{#1}}

\newcommand{\fixme}{}
\newcommand{\incomplete}{\emph{(\ldots incomplete\ldots)}}

\newcommand{\undef}{\emph{undefined}}
\newcommand{\libsie}{\xml{libsie}}


\newcommand{\corever}{1.0}

\title{The \xml{core} SIE Schema, Version \corever{}}

\author{
Brian Downing \\ HBM, Inc. \and
Chris LaReau \\ HBM, Inc.
}

\date{December, 2009}

\begin{document}

\maketitle

\section{Introduction}

SIE provides a very flexible framework for describing almost any kind
of data and associated metadata.  However, some organization and
standardization is neccessary so as to allow the widest amount of
programs to understand the widest amount of data.

SIE tag names can belong to a particular metadata \emph{schema}.  The
schema part of a tag name is the part preceeding a colon.  For
example, tag names beginning with ``\xml{core:}'' belong to the
\xml{core} schema, and will be described in this document.

Not all of the tags below will be present; in fact, none have to be.
For maximum robustness, data reading applications should be able to
work with the absolutely minimum set of tags possible.

In addition to metadata schemas, the \xml{core:schema} tag indicates
which \emph{data schema} is in use for the numerical and binary data
that will be output from a channel.  Referencing the correct data
schema documentation will allow the data to be correctly interpreted.

\section{Metadata schema}

\subsection{\xml{core:description} tag}

The \xml{core:description} tag provides a free-form, natural-language
description of the element it is contained in.  For example:

\begin{siexmlexample}
<tag id="core:description">
  A thermocouple mounted underneath the right front wheel bearing.  
  This channel is gated to only collect data when the temperature is 
  above 60 degrees C.
</tag>
\end{siexmlexample}

\subsection{\xml{core:elapsed\_time} tag}

The \xml{core:elapsed\_time} tag defines the amount of time spent
collecting the data in its container in seconds.

\begin{siexmlexample}
<tag id="core:elapsed_time">160760.4</tag>
\end{siexmlexample}

\subsection{\xml{core:input\_samples} tag}

The \xml{core:input\_samples} tag contains the number of data samples
that were present before any data-reduction algorithms were run.  For
example, in a gated time history, \xml{core:input\_samples} would
contain the number of samples present before gating.

\begin{siexmlexample}
<tag id="core:input_samples">4760</tag>
\end{siexmlexample}

\subsection{\xml{core:label} tag}

The \xml{core:label} tag defines a \emph{label}, which is a short
description intended for placing on a plot, graph, or table.  For a
dimension label, an expected usage would be for an axis label.

\begin{siexmlexample}
<tag id="core:label">Time</tag>
\end{siexmlexample}

\subsection{\xml{core:output\_samples} tag}

The \xml{core:output\_samples} tag contains the number of data samples
that are present after any data-reduction algorithms were run; in
other words, the number of data samples that are actually stored in
the channel.  For example, in a gated time history,
\xml{core:output\_samples} would contain the number of samples stored.

\begin{siexmlexample}
<tag id="core:output_samples">200</tag>
\end{siexmlexample}

\subsection{\xml{core:range\_max} tag}

The \xml{core:range\_max} tag defines the expected maximum bound for
the container dimension.  This can be used, for example, to set plot
range boundries.  It is \emph{not} the maximum data value in the
dimension.

\begin{siexmlexample}
<tag id="core:range_max">1000.0</tag>
\end{siexmlexample}

\subsection{\xml{core:range\_min} tag}

The \xml{core:range\_min} tag is as \xml{core:range\_max}, but defines
the minimum bound for the container dimension.

\begin{siexmlexample}
<tag id="core:range_min">-1000.0</tag>
\end{siexmlexample}

\subsection{\xml{core:sample\_rate} tag}

The \xml{core:sample\_rate} tag defines the sample rate, in Hz, of the
data in its container.

\begin{siexmlexample}
<tag id="core:sample_rate">2500</tag>
\end{siexmlexample}

\subsection{\xml{core:schema} tag}

The \xml{core:schema} tag defines what data schema is in use for this
channel.  The data schema defines how you interpret the data output of
the channel.  For example, if \xml{core:schema} was
\xml{somat:sequential}, you would go to the \xml{somat} schema
document and look up the \xml{sequential} data schema for instructions
on how to interpret it.

\begin{siexmlexample}
<tag id="core:schema">somat:sequential</tag>
\end{siexmlexample}

\subsection{\xml{core:setup\_name} tag}

The \xml{core:setup\_name} tag contains the name of the setup under
which the data in the container has been collected, if applicable.

\begin{siexmlexample}
<tag id="core:setup_name">bearing_temp</tag>
\end{siexmlexample}

\subsection{\xml{core:start\_time} tag}

The \xml{core:start\_time} tag contains the time that data collection
started for the container in ISO 8601 format.

\begin{siexmlexample}
<tag id="core:start_time">2007-01-10T11:49:34-0600</tag>
\end{siexmlexample}

\subsection{\xml{core:stop\_time} tag}

The \xml{core:start\_time} tag is as \xml{core:start\_time}, but
contains the time when data collection stopped.

\begin{siexmlexample}
<tag id="core:stop_time">2007-01-10T12:12:06-0600</tag>
\end{siexmlexample}

\subsection{\xml{core:test\_count} tag}

The \xml{core:test\_count} tag contains the sequence number of the
container test.  For example, on a data acquisition system where a
test can be repeated multiple times, the first test run would have
\xml{core:test\_count} of 1, and the second would be 2.

\begin{siexmlexample}
<tag id="core:test_count">14</tag>
\end{siexmlexample}

\subsection{\xml{core:units} tag}

The \xml{core:schema} tag defines the units of the container dimension
(e.g.\ seconds, millivolts, microstrain, etc.).  For unitless
dimensions (e.g. counts) this tag is absent.

\begin{siexmlexample}
<tag id="core:units">seconds</tag>
\end{siexmlexample}

\subsection{\xml{core:version} tag}

The value of the \xml{core:version} tag defines the version of the
\xml{core} schema in use.  The version described in this document is
\corever{}.

\begin{siexmlexample}
<tag id="core:version">1.0</tag>
\end{siexmlexample}

\section{Data schemas}

Version \corever{} of the \xml{core} schema does not offer any data
schemas.  See the \xml{somat} schema documentation for output schemas
for typical engineering data formats.

\appendix

\section{License}

\begin{verbatim}
Copyright (C) 2005-2010  HBM Inc., SoMat Products

This document is free software; you can redistribute it and/or
modify it under the terms of version 2.1 of the GNU Lesser General
Public License as published by the Free Software Foundation.

This document is distributed in the hope that it will be useful, but
WITHOUT ANY WARRANTY; without even the implied warranty of
MERCHANTABILITY or FITNESS FOR A PARTICULAR PURPOSE.  See the GNU
Lesser General Public License for more details.

You should have received a copy of the GNU Lesser General Public
License along with this document; if not, write to the Free Software
Foundation, Inc., 51 Franklin Street, Fifth Floor, Boston, MA
02110-1301 USA
\end{verbatim}

\end{document}

